\chapter{Considerações Finais}

Este trabalho apresentou o desenvolvimento de uma plataforma didática embarcada para reprodução de registros COMTRADE e emulação de funções de proteção de sistemas elétricos de potência. A solução utiliza duas placas ESP32: a unidade geradora interpreta um registro binário, forma quadros temporizados e reproduz um canal de tensão e um de corrente pelo MCP4728; a unidade receptora adquire esses sinais, estima os fasores fundamentais e avalia as funções ANSI 21, 27, 50, 51, 59 e 67.

A arquitetura alcançada mantém nas ESP32 o caminho de tempo crítico e utiliza o computador para simulação, transferência, configuração e supervisão. O \texttt{config.json} vincula cada COMTRADE às escalas e referências empregadas na recepção, reduzindo ambiguidades entre os domínios elétrico, digital e analógico. Os protocolos de confirmação, os indicadores de qualidade, os eventos seriais e as saídas físicas tornam observável o percurso entre a oscilografia e a decisão de proteção.

Do ponto de vista didático, a plataforma aproxima conceitos de COMTRADE, amostragem, DFT, impedância aparente, pickup, temporização e supervisão direcional de uma experiência reconfigurável em bancada. Sua contribuição principal é a integração funcional e rastreável desses elementos em hardware de baixo custo. O protótipo não substitui uma mala de testes, um relé comercial ou um sistema calibrado, pois não dispõe de condicionamento e isolamento próprios nem foi submetido a uma caracterização metrológica completa.

A observação experimental estabeleceu \SI{1000}{amostras\per\second} como limite operacional para a reprodução coerente dos sinais de \SI{60}{\hertz} avaliados. Acima dessa taxa, houve perda de sincronismo temporal. Esse comportamento, somado às restrições do buffer, da interpolação, do ADC interno e da montagem em protoboard, delimita a validade dos ensaios e evidencia a necessidade de verificar qualidade temporal, saturação e continuidade, e não apenas o estado final de trip.

Como trabalhos futuros, recomendam-se uma matriz quantitativa e repetível de testes para todas as funções implementadas; medição sincronizada dos erros de amplitude, fase e tempo; expansão para três tensões e três correntes; suporte direto a outros formatos e ordenações COMTRADE; condicionamento analógico e isolamento em placa dedicada; avaliação de ADCs externos; curvas inversas para a função 51; comparadores angulares adicionais para a função 67; e comparação experimental com relés comerciais. Essas extensões podem transformar o protótipo funcional em uma bancada didática mais abrangente e metrologicamente caracterizada.
